% =============================================================================
% Chapter 4: Light Research Program
% Input-ready fragment for the lean toy-model research proposal.
% This file intentionally contains no preamble or document environment.
% =============================================================================

\chapter{Light Research Program}
\label{ch:light-research-program}

Light is the first detailed sector in the requirement-discovery program because
it tests the proposed host medium before a particle throat has been solved. A
throat-free ordered slab must already possess the right shear support, guided
mode structure, polarization count, characteristic speed, dispersion, and
isolation from unwanted branches. Those results can be sought directly from the
background brane and its perturbation operator. They therefore expose
incompatible constitutive choices earlier and more cheaply than beginning with
a nonlinear particle core.

The light program contains two nested claims. The first asks whether the common
medium supports ordinary long-wavelength transverse light on the finite slab.
The second asks whether the same medium can also support a finite, freely
propagating classical packet that could serve as an individual photon analog.
Success of the second track requires success of the first, but failure of the
second would not erase a valid conditional light-wave sector. The detailed
sequence of spectral, asymptotic, envelope, and numerical calculations is
maintained in the companion document \emph{Guided Light and the Photon-Soliton
Research Program}. This chapter states the scientific contract and the outputs
that must be handed to the common-medium synthesis.

\section{Research goals and proposed material picture}

The proposed material picture begins with one finite ordered slab embedded in
the de-structured phase of the same medium. The ordered state must support
in-plane shear, whereas the surrounding bulk need not support the same resolved
transverse elastic branch. Light is identified with the lowest physical
transverse branch guided by this finite material structure. It is not a second
substance and no fundamental electromagnetic field is inserted solely to carry
it.

Let the brane possess a local normal axis represented by a unit vector
\(N^A\). The sign of \(N^A\) distinguishes the two throat orientations used by
the charge sector, but ordinary free light should depend at leading order only
on the unsigned axis, through reflection-even structures such as \(N^A N^B\).
For a wave vector tangent to the brane, a light displacement is required to be
both tangent to the slab and transverse to the direction of propagation. In the
four-dimensional parent space, excluding the normal direction and the
propagation direction leaves two independent tangential polarization
directions. The complete guided spectrum must realize that count dynamically;
geometrical counting alone is not sufficient.

Under the currently supplied homogeneous, isotropic quadratic brane action, the
two transverse modes have

\[
\omega_T^2
=
\frac{\mu_R}{\rho_{\mathrm{br}}}k^2,
\qquad
c_\gamma^2
=
\frac{\mu_R}{\rho_{\mathrm{br}}},
\]

where \(\mu_R\) is the effective ordered-state shear modulus and
\(\rho_{\mathrm{br}}\) is the effective inertia density of the brane mode. The
same action contains one in-plane longitudinal branch,

\[
\omega_L^2
=
\frac{B_{\mathrm{comp}}}{\rho_{\mathrm{br}}}k^2,
\qquad
c_L^2
=
\frac{B_{\mathrm{comp}}}{\rho_{\mathrm{br}}}.
\]

At homogeneous quadratic order its transverse and longitudinal blocks decouple. This is a useful
conditional result, but it does not yet establish a stable finite slab, derive
the shear modulus from the parent medium, confine a transverse mode across the
normal direction, or control mixing caused by interfaces, defects,
nonlinearity, and reservoir variables.

Three uses of the transverse sector must remain distinct throughout the
research:

\begin{enumerate}[leftmargin=2.2em,itemsep=0.35em]
  \item \emph{Ordinary light} is a small-amplitude guided transverse excitation
        of the throat-free slab. Plane waves and ordinary spreading packets
        belong to this sector.
  \item A \emph{free photon-packet candidate} is a finite-energy traveling
        object localized across \(w\) and within all three brane directions,
        with only two freely selectable asymptotic polarizations.
  \item A \emph{throat support mode} is a transverse standing mode bound to, or
        resonant around, a nonlinear particle throat. It may exist even if the
        free slab does not support a self-bound photon packet.
\end{enumerate}

The immediate goals are therefore ordered. First, determine whether the
throat-free slab possesses a healthy gapless guided transverse branch with two
polarizations and acceptable leakage. Second, derive its finite-thickness
dispersion and nonlinear couplings from the same common constitutive law. Third,
test whether those outputs permit a stable or sufficiently long-lived
three-dimensional traveling packet. Finally, record every coefficient, profile,
speed, continuum threshold, and unwanted coupling that the light calculation
imposes on gravity, electricity, magnetism, radiation, inertia, and the later
throat-support problem.

\section{Track A: ordinary guided transverse light}

Track A asks for a controlled effective light-wave sector before making the
stronger photon claim. The finite slab must act as a genuine waveguide rather
than merely being declared one. A variable-coefficient transverse reduction has
the schematic form

\[
S_T^{(2)}
=
\frac12\int dt\,d^3x\,dw
\left[
\rho_T(w)|\dot{\mathbf u}_T|^2
-\mu_\parallel(w)|\nabla_\parallel\mathbf u_T|^2
-\mu_w(w)|\partial_w\mathbf u_T|^2
-\mathbf u_T\!\cdot\!\mathbf V_T(w)\mathbf u_T
\right],
\]

with coefficient profiles derived from the solved slab and the selected
shear-supporting constitutive structure. This expression is a diagnostic
template rather than a frozen final action. The actual operator may contain
higher gradients, nonlocal response, coupling to density, order, thickness, or
longitudinal variables, and branch-appropriate damping with its internal
reservoir.

The normal problem must identify a lowest transverse profile \(f_0(w;k)\) and
its projected dispersion \(\omega_0(k)\). On a conservative or inertial branch,
the intended light mode is gapless:

\[
\lambda_0=0,
\qquad
\omega_0(k)\longrightarrow0
\quad\text{as}\quad k\longrightarrow0.
\]

A fixed normal standing frequency that gives
\(\omega^2=\Omega_w^2+c^2k^2\) would instead describe a gapped branch and cannot
serve as ordinary long-wavelength light. On a dissipative or mixed branch, the
corresponding requirement is a brane-coupled physical pole or continuous branch
of the full retarded response whose real frequency approaches zero in the
required way while its damping, spectral weight, and lifetime remain
acceptable.

Calling the bulk ``unable to carry shear'' is not a proof of confinement. A
vanishing bulk shear modulus could yield a bound profile, an evanescent tail, a
floppy continuum, or strong conversion into other material motion. The complete
normal operator and boundary conditions must determine which occurs. A true
bound mode requires a well-defined physical norm, finite energy, and acceptable
normal decay. A leaky mode must instead be treated as an outgoing resonance with
calculated complex frequency, emitted flux, and lifetime. The calculation must
also verify that normalization does not hide a profile whose support spreads
arbitrarily far into the bulk. Useful diagnostics include the branch's
separation from the relevant essential spectrum, its normal localization
length, the fraction of its norm or energy carried by the ordered slab, its
effective three-dimensional inertia under a fixed on-brane amplitude, and the
residue with which a brane observable couples to the mode.

The two transverse polarizations must remain healthy and degenerate on the
homogeneous isotropic background. The full constrained mode census must also
retain the separate longitudinal, normal-interface, order, density, and bulk
branches. Track A does not require those branches to disappear. It requires
that their spectra, couplings, and leakage be known well enough to show that
ordinary transverse propagation is not destabilized or accompanied by
unacceptable conversion, birefringence, attenuation, or additional radiation.
Near a throat or another inhomogeneity, any reactivated
transverse--longitudinal or transverse--bulk mixing becomes a calculable
prediction rather than an omitted effect.

Finite thickness supplies a material length scale and generally modifies the
in-brane dispersion. After projection, the lowest branch may have

\[
\omega_T^2(k)
=
c_\gamma^2k^2
+a_4H_{\mathrm{br}}^2k^4
+a_6H_{\mathrm{br}}^4k^6
+\cdots.
\]

The coefficients \(a_4,a_6,\ldots\) are outputs of the slab reduction. They
must be small enough over the intended ordinary-light regime to preserve the
required approximately nondispersive propagation, yet their signs and
magnitudes may become important for the photon-packet track. They cannot be
selected independently for packet binding after the background dispersion has
been computed.

A successful Track A result therefore consists of more than the homogeneous
formula for \(c_\gamma\). It must deliver a stable slab profile; the complete
transverse eigenfunction; a branch-appropriate gapless classification; two
polarizations; a declared low-dispersion regime; normal localization or an
acceptable resonance lifetime; the surrounding continua and coupling
strengths; and the dependence of these outputs on the same parameters that
enter the force sectors. Until those calculations are completed, the present
light-wave result remains conditional on the assumed quadratic action.

The observer-facing contract should be evaluated over a frequency and
wave-number family, not at one specially favorable point. The phase and group
velocities must be recorded separately, as must any frequency dependence of the
polarization vectors and normal profile. Positive energy and a healthy physical
norm are required throughout the declared light regime. The two polarization
branches should remain degenerate on an isotropic background, and any splitting
introduced by curvature, defects, environmental asymmetry, or nonlinear
intensity must be computed rather than hidden in an averaged speed. Ordinary
free light must also carry no leading reflection-odd electric monopole on a
symmetric background. This requirement does not make the light sector exact
Maxwell theory: the medium may retain a separate longitudinal branch and other
falsifiable material structure. It requires only that the transverse carrier
seen as light be cleanly identifiable and that the additional structure remain
stable, sufficiently weakly coupled, and observationally acceptable in the
stated regime.

Track A must finally provide a normalized carrier for later source calculations.
The same mode norm, energy flux, and polarization basis used to describe free
light must be used when an accelerating oriented throat is tested for emission.
Otherwise the radiation amplitude could be adjusted independently of the
background light sector. This later use does not require a solved throat now,
but it does require Track A to export unambiguous normalization and flux
conventions rather than only a dispersion curve.

\section{Track B: localized classical photon-packet candidate}

A linear guided mode is not yet an individual photon analog. Linear equations
permit plane waves and localized superpositions, but they do not normally
select a finite packet width, prevent diffraction in three dimensions, or
produce a distinguished unit of energy. Guidance across \(w\) solves only the
normal-confinement problem. A packet moving along, for example, the \(z\)
direction may still spread in \(x\), \(y\), and along its envelope even when its
normal profile remains bound.

The preferred candidate is therefore an oscillatory transverse carrier inside
a localized vector envelope,

\[
\mathbf u_T(\mathbf x,w,t)
=
\operatorname{Re}\!\left[
\left(
\mathbf e_1A_1(\mathbf x,t)
+
\mathbf e_2A_2(\mathbf x,t)
\right)
 f_0(w)
 e^{i(k_0z-\omega_0t)}
\right].
\]

The normal profile \(f_0\) is inherited from Track A. The complex envelopes
\(A_1\) and \(A_2\) carry the two independent polarization amplitudes. The
carrier retains a frequency, wavelength, and phase while the envelope localizes
energy, momentum, strain, and stress. Compact support of the displacement itself
is unnecessary; the physical requirement is finite, concentrated energy and
momentum with no unacceptable secular tail.

The proposed binding mechanism combines finite-thickness dispersion with a
reflection-even material dressing excited by the transverse intensity

\[
I=|A_1|^2+|A_2|^2.
\]

The helper is not assumed to be a pure longitudinal displacement. In the finite
slab it may be a derived mixture of in-plane compression, even thickness
change, density response, and order response,

\[
\eta_{\mathrm{helper}}
=
a_L\nabla\!\cdot\!\mathbf u_L
+a_t h_t
+a_n\delta n
+a_\chi\delta\chi_B
+\cdots.
\]

The coefficients and mode composition must come from the complete background
spectrum. At weak amplitude the desired hierarchy is a transverse carrier at
first order, an even helper sourced by intensity at second order, and feedback
on the carrier at third order. If the helper is functionally determined by
\(I\), it is a slaved material dressing rather than a third freely selectable
photon polarization. Its even reflection parity also allows the packet to use
the slab's thickness and compression response without acquiring the leading
odd electric monopole assigned to an oriented throat.

Eliminating the helper through its retarded response may generate a nonlocal,
frequency-dependent self-interaction for the two transverse envelopes. The
general reduction must retain

\[
\eta(\omega,\mathbf q)
=-\alpha\mathcal G_\eta^R(\omega,\mathbf q)I(\omega,\mathbf q).
\]

Only in a controlled quasistatic regime may this reduce schematically to

\[
\eta(\mathbf q)
=-\alpha\mathcal L_\eta^{-1}(\mathbf q)I(\mathbf q),
\]

so that the packet creates a moving guide in the material through which its own
transverse carrier propagates. A focusing response may oppose diffraction; a
defocusing response cannot. A simple local focusing nonlinearity is not enough
if its energy is unbounded below and the packet collapses. Saturation,
competing nonlinearities, higher-gradient terms, or a finite-range nonlocal
helper may be required, but every such stabilizing structure must be part of the
shared medium rather than a photon-only repair.

Two reduced routes are useful. A Korteweg--de Vries-like one-dimensional
solitary shear pulse is an inexpensive test of whether finite-thickness
dispersion and amplitude-dependent propagation can balance at all. It is not
the preferred final photon model because it does not establish lateral
localization, two-polarization carrier structure, or near amplitude-independent
light speed. The principal route is the guided oscillatory carrier with a
three-dimensional vector envelope and slaved helper. A one-dimensional success
is therefore a precursor; it must survive transverse stability, full
three-dimensional continuation, and reconstruction in the parent fields.

The reduced vector system must preserve the symmetries of the parent slab. On a
homogeneous isotropic background, no fixed in-brane direction may be introduced
to collimate one chosen propagation axis. Any effective guide should align with
the packet that creates it. Polarization-dependent nonlinear invariants must be
derived explicitly so that linearly, circularly, and elliptically polarized
families can be compared. A helper sourced only by total intensity may preserve
leading polarization degeneracy, whereas other allowed invariants may split
energies, speeds, widths, or stability. Such splitting is a prediction to be
bounded, not a coefficient to be independently tuned for each polarization.

Existence of one isolated numerical lump would also be insufficient. The desired
result is a connected traveling family over a declared carrier-frequency range,
with calculable relations among amplitude, width, speed, energy, momentum, and
helper strength. Those relations reveal whether localization is robust or
requires tuning to a measure-zero point near collapse, continuum emission, or a
change of spectral branch. They also determine whether ordinary low-intensity
light approaches the linear Track A carrier smoothly. A packet branch that has
no controlled small-amplitude connection, changes speed strongly with energy,
or exists only at a single exceptional frequency may still be an interesting
solitary material excitation, but it would not satisfy the intended broad
classical photon role.

The interpretation of a packet also depends on the selected dynamical branch.
A conservative photon candidate must be a finite-energy traveling periodic
solution with derived energy and momentum, zero net outgoing flux for an exact
bound branch, and appropriate Floquet stability. A conservative resonance must
instead carry a calculated decay rate and lifetime. A dissipative or mixed
candidate must identify the internal reservoir, supplied and dissipated power,
momentum transfer, entropy production, passivity, and depletion or replenishment
within the closed one-medium loop. A freely propagating packet may not draw
unrecorded power from an unchanged bulk.

Within the intended classical stage, a viable packet should propagate at or
very near \(c_\gamma\), have finite energy and momentum, possess no independent
rest configuration with the same conserved data, and approach a light-like
relation such as

\[
E_\gamma\simeq c_\gamma|\mathbf P_\gamma|
\]

in its declared regime. It should remain localized for the required lifetime,
retain exactly two free polarizations, and interact weakly or soliton-like with
other packets in ordinary vacuum conditions. These conditions do not by
themselves derive \(E=\hbar\omega\), photon number, quantum statistics, or
discrete emission and absorption. Those remain outside the present classical
packet test.

\section{Calculations, unwanted channels, and early failure tests}

The light program is organized to reject unsuitable branches before expensive
four-dimensional nonlinear simulation. Its practical sequence is:

\begin{enumerate}[leftmargin=2.2em,itemsep=0.42em]
  \item solve a stable or sufficiently metastable finite slab and freeze the
        branch-appropriate background and reservoir data;
  \item derive the complete constrained linear spectrum, including guided
        transverse, longitudinal, interface, density, order, bulk, and
        reservoir response;
  \item project the lowest transverse branch to obtain \(c_\gamma\), its normal
        profile, finite-thickness dispersion, leakage, and observational range;
  \item derive the first allowed nonlinear invariants and their projections
        onto both transverse polarizations and every candidate helper or loss
        branch;
  \item obtain a vector envelope/helper system from a controlled multiple-scale
        reduction rather than assuming a nonlinear Schrödinger equation;
  \item apply energy, scaling, collapse, continuum, radiation, and transverse
        instability tests before seeking full packet families;
  \item continue surviving reduced solutions to the parent fields, then test
        long-time stability, leakage, and two-packet collisions.
\end{enumerate}

The longitudinal branch is central to the unwanted-channel audit. Its presence
is not itself a failure and it may participate in the even helper. What matters
is whether a moving packet excites a localized comoving response or a persistent
wake. Three speed comparisons must remain distinct: the packet speed
\(v_{\mathrm{packet}}\) relative to the free in-plane longitudinal speed
\(c_L\), the placement of that longitudinal branch relative to the bulk-sound
speed \(c_{s0}\), and the packet speed relative to the bulk continuum. In the
simple homogeneous comparison,

\[
k_w^2
=
k^2\left(\frac{c_L^2}{c_{s0}^2}-1\right),
\]

so \(c_L=c_{s0}\) is a grazing threshold for the free longitudinal branch. It
is not by itself proof of leakage or confinement. Actual radiation also
requires nonzero source projection and interface overlap with an outgoing mode.
The complete finite-slab operator may shift or replace this simple threshold.

The nonlinear helper produces a broader test than the free longitudinal
branch. A traveling packet sources a DC comoving component, carrier harmonics,
sidebands, and possibly additional internal frequencies. Every such component
must be compared with every receiving longitudinal, bulk, thickness, density,
order, odd-electric, and conversion branch using the retarded operator and the
actual source overlap. A DC response may remain localized while a second
harmonic radiates. Conversely, a kinematically open continuum may remain dark
by symmetry or vanishing overlap. Neither outcome may be assumed.

The most useful early gates are summarized in
\cref{tab:light-early-gates}.

\begin{table}[htbp]
\centering
\caption{Early gates for the light program.}
\label{tab:light-early-gates}
\begin{tabularx}{\textwidth}{@{}L{0.25\textwidth}L{0.33\textwidth}Y@{}}
\toprule
\textbf{Gate} & \textbf{Required result} & \textbf{Failure meaning} \\
\midrule
Stable host slab
& Bounded or sufficiently metastable finite ordered slab with positive shear
  support
& The intended light architecture has no viable background. \\
Guided transverse branch
& Branch-appropriate gapless mode, two polarizations, and acceptable normal
  localization or resonance lifetime
& Track A fails for the selected common-medium branch. \\
Ordinary-light regime
& Acceptably small dispersion, attenuation, birefringence, and conversion over
  the declared range
& The mode is not an acceptable ordinary light analog in that regime. \\
Helper sign and boundedness
& Even intensity response focuses when required without an unbounded collapse
  direction
& The proposed self-guiding mechanism fails or requires a different common
  constitutive branch. \\
Radiation audit
& DC, harmonic, and sideband coupling to all unwanted branches is absent,
  suppressed, or quantitatively acceptable
& The packet produces an unacceptable wake, leakage, drag, or lifetime. \\
Three-dimensional existence
& Finite-energy localized traveling family survives transverse and parent-field
  continuation
& Track B fails even if Track A remains conditionally viable. \\
Stability and collisions
& No growing physical mode and weak or soliton-like ordinary interactions
& The candidate is not a persistent free photon analog. \\
Cross-sector inheritance
& Every useful coefficient and profile remains compatible with the force and
  throat sectors without retuning
& A light-only success does not define one common medium. \\
\bottomrule
\end{tabularx}
\end{table}

These gates deliberately distinguish levels of failure. Absence of a stable
host slab or a healthy guided transverse branch rejects the selected ordinary
light architecture. Failure of a particular helper mode rejects that binding
mechanism, not necessarily all possible common-medium branches. Failure of all
permitted shared self-binding mechanisms would leave classical guided light but
not the intended localized free-photon ontology. A packet found only after
adding a private coefficient or ignoring its force-sector consequences does not
count as a success.

Several of the early tests can produce useful negative results without a full
packet simulation. A virial or scaling identity may show that all local focusing
branches collapse or disperse. The projected helper Green function may have the
wrong sign. A required harmonic may intersect an outgoing continuum with
nonzero overlap for every admissible speed. The one-dimensional precursor may
possess an unavoidable transverse instability. A non-normal retarded operator
may make an apparently narrow resonance extremely sensitive to perturbations.
Each result should be recorded at the narrowest justified level: it may reject a
coefficient sign, a helper eigenmode, a dynamical branch, or the complete packet
track. These analytic and reduced-model gates are valuable precisely because
they prevent a difficult numerical search from being interpreted as evidence
for a solution when the governing structure already forbids one.

\section{Requirements handed to the common-medium synthesis}

The light program must return normalized requirements and calculated outputs,
not merely the statement that ``the brane carries waves.'' The principal
handoff is summarized in \cref{tab:light-handoff}.

\begin{table}[htbp]
\centering
\caption{Principal light-sector handoff to the common-medium synthesis.}
\label{tab:light-handoff}
\begin{tabularx}{\textwidth}{@{}L{0.27\textwidth}X@{}}
\toprule
\textbf{Category} & \textbf{Required output or restriction} \\
\midrule
Background material
& Stable finite slab profile, selected thickness, positive ordered-state shear
  support, and explicit conservative, relaxational, or mixed branch with its
  reservoir closure. \\
Guided carrier
& Lowest physical transverse profile \(f_0(w;k)\), two polarization vectors,
  gapless classification, \(c_\gamma\), mode norm, brane participation, and
  leakage or resonance lifetime. \\
Dispersion and scales
& Valid ordinary-light range; coefficients such as \(a_4,a_6,\ldots\); relevant
  continuum thresholds; and any relation among \(H_{\mathrm{br}}\), localization
  length, and carrier wavelength. \\
Complete spectrum
& Longitudinal, interface, density, order, bulk, and reservoir branches with
  their parity, characteristic speeds, damping, source overlaps, and dangerous
  low-threshold channels. \\
Nonlinear response
& Allowed transverse invariants; helper eigenvector and parity; focusing or
  defocusing sign; saturation, nonlocality, memory, and boundedness; and the
  same coefficients' effects on the force sectors. \\
Radiative closure
& Derived power or leakage for packet motion, together with the normalized
  outgoing modes, projectors, and loss channels needed to evaluate later
  accelerating-throat emission into both transverse polarizations and every
  odd, even, longitudinal, bulk, and conversion branch. \\
Packet outcome
& Either a stable or sufficiently long-lived three-dimensional family with its
  energy, momentum, speed, polarization, lifetime, and collision behavior, or a
  bounded negative result identifying which requirement cannot be met. \\
Frozen inheritance
& No photon-specific terms; every profile, coefficient, speed, branch choice,
  and reservoir datum is exported to Candidate Parent System V1 and reopened if
  an upstream choice changes. \\
\bottomrule
\end{tabularx}
\end{table}

Several cross-sector burdens follow immediately. The same slab thickness and
constitutive coefficients that guide light also determine the normal parity
spectrum used by electricity, the even thickness response used in finite-size
forces, the longitudinal and bulk modes that control drag and extra radiation,
and the transverse operator later used by throat support and electromagnetic
emission. A gapped even helper may be useful for packet localization and for
screening unwanted thickness forces, but its speed, range, and nonlinear
couplings must satisfy both roles. Finite-thickness dispersion useful for
binding cannot produce unacceptable vacuum dispersion. A light branch close to
a longitudinal or bulk continuum may constrain the characteristic-speed region
available to gravity and the electric sector. Most decisively, the same
transverse coupling structure must later accept energy and momentum from an
accelerating oriented throat in both polarizations without opening excessive
competing channels.

The common-medium synthesis in \cref{ch:cross-sector-integration} will therefore
not score light with a single pass or fail label. It will intersect a light
requirement region containing the stable-slab, guided-mode, ordinary-propagation,
helper, packet, leakage, and branch-closure conditions with the corresponding
regions from the far-field force program. A robust common region would justify
freezing those light outputs in Candidate Parent System V1 and rederiving them
from its complete constrained spectrum. A fragile region near continuum,
instability, or collapse thresholds would remain a risk even if one numerical
packet were found. An empty intersection would reject the candidate combination
rather than invite sector-specific retuning.

Light thus performs two roles in the proposal. Track A determines whether the
medium can carry ordinary two-polarization guided radiation at all. Track B
asks whether the same material can turn that radiation into a localized
classical photon analog without sacrificing ordinary vacuum behavior. Together
they provide an early, unusually direct test of the common background and a
precise set of inherited quantities for the force, throat, and reservoir work
that follows.
